% Variants of Newtons Method for Obtaining the Apsis of Collision

\subsection{Iterating towards the Exact Solution}
%%%%%%%%%%%%%%%%%%%%%%%%%%%%%%%%%%%%%%%%%%%%%%%%%
\label{s:sn_apsis_iterating}
%
We consider three options for iteration:
%
\subsubsection{Newton’s method}
%
Newton's method to solve \cref{eq:sn_apsis_eq} for $R = R_0$ is given by \cref{eq:sn_apsis_newton}.
The initial condition \cref{eq:sn_apsis_ic} guarantees convergence of the iterations for screening functions with $\Phi'(R) < 0$ and $\Phi''(r) > 0$.

\subsubsection{Modified Newton's method}
%
Newton's method applied to \cref{eq:sn_apsis_eq} can also be interpreted as constructing the tangent lines to the terms $\Phi(R)/\varepsilon$ and $R - P^2/R$ and intersecting the two tangents, see the dashed lines in \cref{fig:apsis-iteration}.
%
\begin{figure}[htbp]
    \centering
    \includegraphics[scale=0.6]{apsis_iteration.pdf}
    \caption{Illustration of Newton’s method (linearization of both \(\Phi(R)\) and \(R-P^2/R\)) and the method based on linearization of \(\Phi(R)\) only.
    The latter comes closer to the actual solution \(R_0\).
    For both methods \(R_0^{(i)} < R_0^{(i+1)} < R_0\) provided \(R_0^{(i)} < R_0\).}
    \label{fig:apsis-iteration}
\end{figure}

Instead, one could also intersect the tangent to $\Phi(R)/\varepsilon$ with the function $R - P^2/R$.
The resulting quadratic equation has the solution%
%
\footnote{The other sign in front of the square root is not possible since all coefficients $a$, $b$, and $c$ are positive, which would result in a negative value of $R_0^{(i+1)}$.
}
%
\begin{equation}
    \label{eq:sn_apsis_linphi}
    R_0^{(i+1)} = \frac{ b + \sqrt{ b^2 + 4ac } }{ 2a }
\end{equation}
%
with
%
\begin{equation}
    \label{eq:sn_apsis_linphi_params}
    a = \varepsilon-\Phi'(R_0^{(i)}), \quad
    b = \Phi(R_0^{(i)}) - R_0^{(i)} \Phi'(R_0^{(i)}), \quad
    c = \varepsilon P^2
\end{equation}
%
As shown in \cref{fig:apsis-iteration}, both methods warrant \(R_0^{(i)} < R_0^{(i+1)} < R_0\) provided \(R_0^{(i)} < R_0\), as long as $\Phi''(R) > 0$.
In addition, \cref{fig:apsis-iteration} demonstrates that \cref{eq:sn_apsis_linphi} is a better approximation to \(R_0\) than \cref{eq:sn_apsis_newton}. This comes at the expense of another square root evaluation.

\subsubsection{Newton's method applied to $\log(f(R))$}
%
Instead of using \cref{eq:sn_apsis_newton} with \(f(R)\) from \cref{eq:sn_apsis_eq}, one can move the first and third term of $f(R)$  to the right-hand side of $f(R) = 0$, take the logarithm, and move the right-hand side back to the left-hand side.
One then obtains

\begin{equation}
    \label{eq:sn_apsis_newton_log}
    f_{\log}(R) = \log \left[ \frac{ 1 }{ \varepsilon } \Phi(R) \left/
        \left( R - \frac{ P^2 }{ R } \right)
        \right. \right]
        = 0
\end{equation}
%
The function \(f_{\log}(R)\) can be used with Newton’s method \cref{eq:sn_apsis_newton}.
It maintains the properties \(f_{\log}'(R) < 0\) and \(f_{\log}''(R) > 0\) for screening functions that are given by a sum of exponentials.
Therefore, Newton’s method is guaranteed to converge if one starts with an initial condition to the left of the solution.
The idea of taking the logarithm is to reduce the nonlinearity of \(f(R)\), but obviously comes at the expense of the evaluation of another special function.

\subsection{Convergence}
%%%%%%%%%%%%%%%%%%%%%%%%
\label{s:sn_apsis_performance}
%
\Cref{fig:iter-apsis-legacy} evaluates the perfomance of the combinations of initial condition and iteration algorithm for the ZBL screening function.
It shows the number of iterations required to reach a relative accuracy of \(10^{-3}\) between the last two calculated values of \(R_0\).
Usually this means that the last value is accurate to \(10^{-6}\), assuming quadratic convergence has been reached.

\begin{figure}[htbp]
    \centering
    \includegraphics[scale=0.8]{iter_apsis_zbl.pdf}
    \caption{Number of iterations required to reach a relative accuracy of \(10^{-3}\) in \(R_0\) as a function of energy and impact parameter.
    Left column: using \(R_0^{(0)} = P\) \cref{eq:sn_apsis_ic_p} as initial condition.
    Right column: using \cref{eq:sn_apsis_ic} as initial condition.
    The columns differ by the iteration method used:
    First row: Newton.
    Second row: \cref{eq:sn_apsis_linphi} in the first iteration, then Newton.
    Third row: \cref{eq:sn_apsis_linphi} for all iterations.
    Fourth row: \cref{eq:sn_apsis_linphi} in the first iteration, then logarithmic Newton \cref{eq:sn_apsis_newton_log}.
    The interatomic potential used was the universal ZBL potential.}
    \label{fig:iter-apsis-legacy}
\end{figure}

In the left column, results for the initial condition \cref{eq:sn_apsis_ic_p} (\(R_0^{(0)} = P\)) are shown.
As expected, fewer iterations are required for large energies and large impact parameters (small deflection angles) than for small energies and small impact parameters (large deflection angles), since \(R_0\) is closer to \(P\).
For large deflection angles, more than 10 iterations may be required using ordinary Newton or \cref{eq:sn_apsis_linphi}.
Significant improvement is obtained by using the logarithmic function \cref{eq:sn_apsis_newton_log}, limiting the number of iterations to 6.
Note that Newton’s method with the logarithmic function cannot be started with \(R_0^{(0)} = P\), since this would lead to \(f_{\log}(R \rightarrow P) \rightarrow \infty\).
Therefore, \cref{eq:sn_apsis_linphi} has been used in the first iteration.

In the right column, the table of \(R_0^{(0,\varepsilon)}(\varepsilon)\) for head-on collisions has been used as described above.
Compared to the left column, the number of iterations at low energies and small impact parameters is significantly reduced (not couting the iterations necessary to obtain the table).
There are only slight differences in the iteration counts between the different iteration methods.
Using \cref{eq:sn_apsis_linphi} and \cref{eq:sn_apsis_newton_log} save one iteration in rare cases (compare first, second, and fourth row), while using \cref{eq:sn_apsis_linphi} beyond the first collision has no benefit (compare second and third row).
Since convergence is only marginally improved by using \cref{eq:sn_apsis_linphi} and \cref{eq:sn_apsis_newton_log}, while additional evaluations of a square root or a logarithm is necessary, using ordinary Newton’s method seems to be sufficient.

In summary, when no table of \(R_0^{(0,\varepsilon)}(\varepsilon)\) for head-on collisions is used, it is recommended to use \cref{eq:sn_apsis_linphi} in the first iteration and Newton’s method with the logarithmic function \cref{eq:sn_apsis_newton_log} for the remaining iterations. Significant acceleration is obtained by using the table of \(R_0^{(0,\varepsilon)}(\varepsilon)\) for head-on collisions to obtain an initial condition.
Then it is sufficient to use ordinary Newton’s method.
Note that the table needs to be calculated only once for a given interatomic potential.
